\subsubsection{Characteristic \slash{} Riemann-Invariant Boundary Condition}
\label{SEC:BC:Farfield:Characteristic}

The characteristic boundary condition is an inflow\slash outflow boundary
condition that works for both subsonic and supersonic flows.
%
This boundary condition uses the Riemann invariants to compute the right state
variables based on what flow regime the boundary condition state is in, which is
determined by the magnitude and direction normal to the boundary face of the
interior velocity.
%
Here are a few comments\slash warnings about using this boundary condition
before continuing on with describing how it is formulated:
%
\begin{itemize}[itemsep=5pt,topsep=5pt,parsep=0pt]
  \item When used as a subsonic inflow boundary condition, there is no
    guarantee that the total conditions at the inflow will be enforced.
    This should not be a problem if the inflow is significantly far
    upstream from the computational body of interest (\forexample external
    flows).
  \item This works well as a farfield boundary condition and has not shown any
    problems when the boundary condition velocity is aligned tangent to the
    boundary face.
  \item This works well for an outflow boundary condition for both internal and
    external flows, whether they are subsonic, supersonic, or mixed.
  \item If used for an outflow boundary condition, the prescribed velocity for
    the boundary condition is important for subsonic outflows, whereas the
    boundary condition velocity is ignored for the subsonic outflow boundary
    condition in \sref{SEC:BC:Outflow:Subsonic}. It should be noted that,
    similar to when used for subsonic inflows, this boundary condition applied
    to subsonic outflows does not guarantee that any specified boundary
    condition state will be held constant.
\end{itemize}

%%%%%%%%%%%%%%%%%%%%%%%%%%%%%%%%%%%%%%%%%%%%%%%%%%%%%%%%%%%%%%%%%%%%%%%%%%%%%%%%

Recalling that \normal{} is the outward-facing unit normal for the boundary
face, let $v_{n,\subL}=\vec{u}_\subL\cdot\normal$ be the magnitude of the
interior velocity normal to the boundary face, and
$v_{n,\subBC}=\vec{u}_\subBC\cdot\normal$ be the same but for the prescribed
boundary condition velocity.
%
Similarly, let $a_\subL$ and $a_\subBC$ respectively be the speed of sound for
the interior and boundary condition states.
%
In order to compute the final right state, we need to compute four quantities:
$R^+$, the outgoing Riemann invariant; $R^-$, the incoming Riemann invariant;
$\vec{v}_{t,\subR}$, the tangential velocity vector at the boundary
face; and $s_\subR$, the entropy for the right state.

%%%%%%%%%%%%%%%%%%%%%%%%%%%%%%%%%%%%%%%%%%%%%%%%%%%%%%%%%%%%%%%%%%%%%%%%%%%%%%%%

The outgoing Riemann invariant, $R^+$, is evaluated as
%
\[
  R^+ =
  \begin{cases}
    v_{n,\subBC} + \dfrac{2 a_\subBC}{\gamma-1} &
      \text{\;if } \dfrac{v_{n,\subBC}}{a_\subBC} \leq -1
    \quad \text{(supersonic inflow)} \\[2.5ex]
    v_{n,\subL} + \dfrac{2 a_\subL}{\gamma-1} &
      \text{\;otherwise}
  \end{cases}
\]
The incoming Riemann invariant, $R^-$, is evaluated as
%
\[
  R^- =
  \begin{cases}
    v_{n,\subL} - \dfrac{2 a_\subL}{\gamma-1} &
      \text{\;if } \dfrac{v_{n,\subL}}{a_\subL} \geq 1
    \quad \text{(supersonic outflow)} \\[2.5ex]
    v_{n,\subBC} - \dfrac{2 a_\subBC}{\gamma-1} &
      \text{\;otherwise}
  \end{cases}
\]
%
The tangential velocity vector, $\vec{v}_{t,\subR}$, is evaluated as
%
\[
  \vec{v}_{t,\subR} =
  \begin{cases}
    \vec{u}_\subBC - \left(\vec{u}_\subBC\cdot\normal\right)\normal &
      \text{\;if } v_{n,\subL} \leq 0 \quad \text{(inflow)} \\[1.5ex]
    \vec{u}_\subL - \left(\vec{u}_\subL\cdot\normal\right)\normal &
      \text{\;if } v_{n,\subL} > 0 \quad \text{(outflow)}
  \end{cases}
\]
%
Finally, the entropy, $s_\subR$, is evaluated as
\[
  s_\subR =
  \begin{cases}
    \dfrac{p_\subBC}{\rho^{\gamma}_{\subBC}} &
      \text{\;if } v_{n,\subL} \leq 0 \quad \text{(inflow)} \\[2.5ex]
    \dfrac{p_\subL}{\rho^{\gamma}_{\subL}} &
      \text{\;if } v_{n,\subL} > 0 \quad \text{(outflow)}
   %  \\[2.0ex]
   %p_\subBC / \rho^{\gamma}_{\subBC} &
   %  \text{\;if } v_{n,\subL} \leq 0 \quad \text{(inflow)} \\[1.0ex]
   %p_\subL / \left(\rho^{\gamma}_{\subL}\right) &
   %  \text{\;if } v_{n,\subL} > 0 \quad \text{(outflow)}
  \end{cases}
\]

%%%%%%%%%%%%%%%%%%%%%%%%%%%%%%%%%%%%%%%%%%%%%%%%%%%%%%%%%%%%%%%%%%%%%%%%%%%%%%%%

These four quantities are used to compute the right state in the following way.
%
First, the characteristic normal velocity is computed by averaging the Riemann
invariants
%
\[
  v_{n,\subR} = \frac{1}{2}\left(R^+ + R^-\right)
\]
%
From there, the right state velocity can be determined by simply adding the
normal and tangential components.
%
\[
  \vec{u}_\subR = v_{n,\subR}\normal + \vec{v}_{t,\subR}
\]
%
Next, the speed of sound is computed.
%
\[
  a_\subR = \frac{\gamma-1}{4} \left( R^+ - R^- \right)
\]
%
Finally, the right state density and pressure can be computed as
%
\begin{alignat*}{1}
  \rho_\subR &
      = \left(\frac{a^2_\subR}{\gamma s_\subR}\right)^{\tfrac{1}{\gamma-1}}
      \\[1.0ex]
  p_\subR &
      = \frac{a^2_\subR \rho_\subR}{\gamma}
\end{alignat*}

%%%%%%%%%%%%%%%%%%%%%%%%%%%%%%%%%%%%%%%%%%%%%%%%%%%%%%%%%%%%%%%%%%%%%%%%%%%%%%%%

\devNotes{Currently, the characteristic speed of sound is calculated as $a_c =
(R^+ - R^-)(\gamma-1)/4$. From the first paragraph on page 379 of Hirsch vol.
2, this does not enforce total temperature, and that the equation $a_c^2 =
(H_\subBC - \vec{u}^2_\subR)(\gamma-1)/2$ should be used instead. Also the way
pressure is computed should be changed from $p=s\rho^{\gamma}$ to $p=a_c^2
\rho/\gamma$. These are equivalent, but the latter is just nicer without the
$\gamma$ in the exponent.
%
When switching to this notation, need to clarify what the subscript $c$
means in $a_c$ and how it relates to $a_\subR$.
}

%###############################################################################

